\documentclass[12pt, a4paper]{article}



\usepackage{fontspec}
\usepackage{polyglossia}
\usepackage{csquotes}

\setmainlanguage{russian}
\setotherlanguages{english}

% download "Linux Libertine" OTF-fonts:
% http://www.linuxlibertine.org/index.php?id=91&L=1
\setmainfont{Linux Libertine O} % or Helvetica, Arial, Cambria
% why do we need \newfontfamily:
% http://tex.stackexchange.com/questions/91507/
\newfontfamily{\cyrillicfonttt}{Linux Libertine O}
\newfontfamily{\cyrillicfont}{Linux Libertine O}
\newfontfamily{\cyrillicfontsf}{Linux Libertine O}

\usepackage{etoolbox} % to use ifdef, must be after babel


\usepackage[paper=a4paper, top=13.5mm, bottom=13.5mm, left=16.5mm, right=13.5mm, includefoot]{geometry}

\usepackage{etex} % расширение классического tex
% в частности позволяет подгружать гораздо больше пакетов, чем мы и займёмся далее

\usepackage{floatrow} % сильно вниз если сдвинуть - ругается!


\usepackage{makeidx} % для создания предметных указателей
\usepackage{verbatim} % для многострочных комментариев
%\usepackage[pdftex]{graphicx} % для вставки графики
% omit pdftex option if not using pdflatex


%\usepackage{dsfont} % шрифт для единички с двойной палочкой (для индикатора события)
\usepackage{bbm} % шрифт - двойные буквы


\usepackage[usenames, dvipsnames, svgnames, table, rgb]{xcolor}

\usepackage{colortbl}

\usepackage{comment} % для комментирования кусков

% пакет для тестов:
\usepackage[box, % запрет на перенос вопросов
nopage, % убираем колонтитулы страницы
insidebox, % ставим буквы в квадратики
separateanswersheet, % добавляем бланк ответов
nowatermark, % отсутствие надписи "Черновик"
indivanswers,  % показываем верные ответы
%answers,
lang=RU, % локализация слов "нет верных ответов", "вопрос" и тд
completemulti % добавлять "нет правильного ответа" во всех вопросах множественного выбора
]{automultiplechoice}


\usepackage{multicol}
\usepackage{multirow} % Слияние строк в таблице


\usepackage[colorlinks, hyperindex, unicode, breaklinks]{hyperref} % гиперссылки в pdf

\usepackage{dcolumn}

\usepackage{amssymb}
\usepackage{amsmath}
\usepackage{amsthm}
\usepackage{epsfig}
\usepackage{bm}
\usepackage{color}

\usepackage{todonotes} % для вставки в документ заметок о том, что осталось сделать
% \todo{Здесь надо коэффициенты исправить}
% \missingfigure{Здесь будет картина Последний день Помпеи}
% команда \listoftodos — печатает все поставленные \todo'шки


\usepackage{textcomp}  % Чтобы в формулах можно было русские буквы писать через \text{}

%\usepackage{embedfile} % Чтобы код LaTeXа включился как приложение в PDF-файл

\usepackage{subfigure} % для создания нескольких рисунков внутри одного

\usepackage{tikz, pgfplots} % язык для рисования графики из latex'a
\usetikzlibrary{trees} % прибамбас в нем для рисовки деревьев
\usetikzlibrary{arrows} % прибамбас в нем для рисовки стрелочек подлиннее
\usepackage{tikz-qtree} % прибамбас в нем для рисовки деревьев




\usepackage{enumitem}


%\embedfile[desc={Исходный LaTeX файл}]{\jobname.tex} % Включение кода в выходной файл
%\embedfile[desc={Стилевой файл}]{title_bor_utf8.tex}



% вместо горизонтальной делаем косую черточку в нестрогих неравенствах
\renewcommand{\le}{\leqslant}
\renewcommand{\ge}{\geqslant}
\renewcommand{\leq}{\leqslant}
\renewcommand{\geq}{\geqslant}

% делаем короче интервал в списках
\setlength{\itemsep}{0pt}
\setlength{\parskip}{0pt}
\setlength{\parsep}{0pt}

% свешиваем пунктуацию (т.е. знаки пунктуации могут вылезать за правую границу текста, при этом текст выглядит ровнее)
\usepackage{microtype}

% более красивые таблицы
\usepackage{booktabs}
% заповеди из докупентации:
% 1. Не используйте вертикальные линни
% 2. Не используйте двойные линии
% 3. Единицы измерения - в шапку таблицы
% 4. Не сокращайте .1 вместо 0.1
% 5. Повторяющееся значение повторяйте, а не говорите "то же"


% DEFS
\def \mbf{\mathbf}
\def \msf{\mathsf}
\def \mbb{\mathbb}
\def \tbf{\textbf}
\def \tsf{\textsf}
\def \ttt{\texttt}
\def \tbb{\textbb}

\def \wh{\widehat}
\def \wt{\widetilde}
\def \ni{\noindent}
\def \ol{\overline}
\def \cd{\cdot}
\def \bl{\bigl}
\def \br{\bigr}
\def \Bl{\Bigl}
\def \Br{\Bigr}
\def \fr{\frac}
\def \bs{\backslash}
\def \lims{\limits}
\def \arg{{\operatorname{arg}}}
\def \dist{{\operatorname{dist}}}
\def \VC{{\operatorname{VCdim}}}
\def \card{{\operatorname{card}}}
\def \sgn{{\operatorname{sign}\,}}
\def \sign{{\operatorname{sign}\,}}
\def \xfs{(x_1,\ldots,x_{n-1})}
\def \Tr{{\operatorname{\mbf{Tr}}}}
\DeclareMathOperator*{\argmin}{arg\,min}
\DeclareMathOperator*{\argmax}{arg\,max}
\DeclareMathOperator*{\amn}{arg\,min}
\DeclareMathOperator*{\amx}{arg\,max}
\def \cov{{\operatorname{Cov}}}
\DeclareMathOperator{\Var}{Var}
\DeclareMathOperator{\Cov}{Cov}
\DeclareMathOperator{\Corr}{Corr}
\DeclareMathOperator{\grad}{grad}
\DeclareMathOperator{\trace}{trace}
\DeclareMathOperator{\rank}{rank}

\def \xfs{(x_1,\ldots,x_{n-1})}
\def \ti{\tilde}
\def \wti{\widetilde}


\def \mL{\mathcal{L}}
\def \mW{\mathcal{W}}
\def \mH{\mathcal{H}}
\def \mC{\mathcal{C}}
\def \mE{\mathcal{E}}
\def \mN{\mathcal{N}}
\def \mA{\mathcal{A}}
\def \mB{\mathcal{B}}
\def \mU{\mathcal{U}}
\def \mV{\mathcal{V}}
\def \mF{\mathcal{F}}

\def \R{\mbb R}
\def \N{\mbb N}
\def \Z{\mbb Z}
\def \P{\mbb{P}}
%\def \p{\mbb{P}}
\def \E{\mbb{E}}
\def \D{\msf{D}}
\def \I{\mbf{I}}

\def \a{\alpha}
\def \b{\beta}
\def \t{\tau}
\def \dt{\delta}
\def \e{\varepsilon}
\def \ga{\gamma}
\def \kp{\varkappa}
\def \la{\lambda}
\def \sg{\sigma}
\def \sgm{\sigma}
\def \tt{\theta}
\def \ve{\varepsilon}
\def \Dt{\Delta}
\def \La{\Lambda}
\def \Sgm{\Sigma}
\def \Sg{\Sigma}
\def \Tt{\Theta}
\def \Om{\Omega}
\def \om{\omega}


\def \ni{\noindent}
\def \lq{\glqq}
\def \rq{\grqq}
\def \lbr{\linebreak}
\def \vsi{\vspace{0.1cm}}
\def \vsii{\vspace{0.2cm}}
\def \vsiii{\vspace{0.3cm}}
\def \vsiv{\vspace{0.4cm}}
\def \vsv{\vspace{0.5cm}}
\def \vsvi{\vspace{0.6cm}}
\def \vsvii{\vspace{0.7cm}}
\def \vsviii{\vspace{0.8cm}}
\def \vsix{\vspace{0.9cm}}
\def \VSI{\vspace{1cm}}
\def \VSII{\vspace{2cm}}
\def \VSIII{\vspace{3cm}}



\newcommand{\dx}[1]{\,\mathrm{d}#1} % для интеграла: маленький отступ и прямая d
\newcommand{\ind}[1]{\mathbbm{1}_{\{#1\}}} % Индикатор события
%\renewcommand{\to}{\rightarrow}
\newcommand{\eqdef}{\mathrel{\stackrel{\rm def}=}}
\newcommand{\iid}{\mathrel{\stackrel{\rm i.\,i.\,d.}\sim}}
\newcommand{\const}{\mathrm{const}}

%на всякий случай пока есть
%теоремы без нумерации и имени
\newtheorem*{theorem}{Теорема}

%"Определения","Замечания"
%и "Гипотезы" не нумеруются
\newtheorem*{definition}{Определение}
%\newtheorem*{rem}{Замечание}
%\newtheorem*{conj}{Гипотеза}

%"Теоремы" и "Леммы" нумеруются
%по главам и согласованно м/у собой
%\newtheorem{theorem}{Теорема}
%\newtheorem{lemma}[theorem]{Лемма}

% Утверждения нумеруются по главам
% независимо от Лемм и Теорем
%\newtheorem{prop}{Утверждение}
%\newtheorem{cor}{Следствие}
 % use local copy

\def \useR{$[$R$]$ }

%% эконометрические сокращения
\def \hb{\hat{\beta}}
\DeclareMathOperator{\sVar}{sVar}
\DeclareMathOperator{\sCov}{sCov}
\DeclareMathOperator{\sCorr}{sCorr}


\def \hs{\hat{s}}
\def \hy{\hat{y}}
\def \hY{\hat{Y}}
\def \he{\hat{\varepsilon}}
\def \v1{\vec{1}}
\def \cN{\mathcal{N}}
\def \e{\varepsilon}
\def \z{z}

\def \hVar{\widehat{\Var}}
\def \hCorr{\widehat{\Corr}}
\def \hCov{\widehat{\Cov}}

\DeclareMathOperator{\tr}{tr}
\DeclareMathOperator*{\plim}{plim}

%% лаг
\renewcommand{\L}{\mathrm{L}}

%% алая и белая розы
%% запускается так: \WhiteRose[масштаб], например, \WhiteRose[0.5]
\newcommand{\WhiteRose}[1]{\begingroup
\setbox0=\hbox{\includegraphics[scale=#1]{figures/Yorkshire_rose.pdf}}%
\parbox{\wd0}{\box0}\endgroup}

\newcommand{\RedRose}[1]{\begingroup
\setbox0=\hbox{\includegraphics[scale=#1]{figures/Lancashire_rose.pdf}}%
\parbox{\wd0}{\box0}\endgroup}

\newcommand{\WhiteRoseLine}{
\begin{center}
\WhiteRose{0.3} Версия Белой Розы \WhiteRose{0.3}
\end{center}}

\newcommand{\RedRoseLine}{
\begin{center}
\RedRose{0.3} Версия Алой Розы \RedRose{0.3}
\end{center}}
 % use local copy

% \usepackage{minted}

\unitlength=0.6mm


%%%%%%%%%%%%%%%%%% вставки
%%%%%%%%%%%%%%%%%%%%%%%%%%%%%%%%%%%%%%% Списки без уродских отступов
\newenvironment{enumerate*}{
\begin{enumerate}
  \setlength{\itemsep}{0pt}
  \setlength{\parskip}{0pt}
  \setlength{\parsep}{0pt}
}{\end{enumerate}}

\newenvironment{itemize*}{
\begin{itemize}
  \setlength{\itemsep}{0pt}
  \setlength{\parskip}{0pt}
  \setlength{\parsep}{0pt}
}{\end{itemize}}

\abovedisplayskip=0mm
\abovedisplayshortskip=0mm
\belowdisplayskip=0mm
\belowdisplayshortskip=0mm
%%%%%%%%%%%%%%%%%%%%%%%%%%%%%%%%%%%%%%%%%%%%%%%%%%%%%%%%%%%%%%%%%%%%%%

\newenvironment{centered}{%
  \begin{list}{}{%
    \topsep0pt
  }
  \centering
  \item[]
}
{\end{list}}
%%%%%%%%%%%%%%%%%%%%%%%%%%%%%%%%%%%%%%%%%%%%%%%%%%%%%%%%%%%%%%%%%%%%%%%%%

\usepackage[sorting=none, backend=biber]{biblatex}


\addbibresource{metrics_hse_exams.bib}


\AddEnumerateCounter{\asbuk}{\russian@alph}{щ} % для списков с русскими буквами
\setlist[enumerate, 2]{label=\asbuk*),ref=\asbuk*} % списки уровня 2 будут буквами а) б) ...



\begin{document}




\element{2017_fall_retake_1}{ % в фигурных скобках название группы вопросов
%  %\AMCnoCompleteMulti
\begin{questionmult}{1} % тип вопроса (questionmult — множественный выбор) и в фигурных — номер вопроса

Для набора панельных данных истинна спецификация модели со случайными эффектами, однако Вовочка оценивает модель с фиксированными эффектами. Вовочкины ценки коэффициентов $\beta$ окажутся

\begin{multicols}{2} % располагаем ответы в 3 колонки
\begin{choices} % опция [o] не рандомизирует порядок ответов
       \correctchoice{состоятельными и неэффективными}
       \wrongchoice{несостоятельными}
       \wrongchoice{состоятельными и эффективными}
       \wrongchoice{смещёнными и неэффективными}
       \wrongchoice{несмещёнными и эффективными}
    \end{choices}
   \end{multicols}
\end{questionmult}
}


\element{2017_fall_retake_1}{ % в фигурных скобках название группы вопросов
%  %\AMCnoCompleteMulti
\begin{questionmult}{2} % тип вопроса (questionmult — множественный выбор) и в фигурных — номер вопроса

Винни-Пух пытается понять, от каких переменных может зависеть его потребление мёда. Собрав 100 разных переменных, он построил 100 парных регрессий и проверил в них значимость коэффициента при каждой из переменных на уровне значимости 0.05. Пятачок понимает, что все 100 переменных не имеют никакого отношения к потреблению мёда и на самом деле просто случайные числа. Помогите Пятачку определить, сколько значимых переменных скорее всего найдёт Винни-Пух.

\begin{multicols}{3} % располагаем ответы в 3 колонки
\begin{choices} % опция [o] не рандомизирует порядок ответов
       \correctchoice{5}
       \wrongchoice{10}
       \wrongchoice{0}
       \wrongchoice{100}
       \wrongchoice{Не хватает данных для ответа}
    \end{choices}
   \end{multicols}
\end{questionmult}
}


\element{2017_fall_retake_1}{ % в фигурных скобках название группы вопросов
%  %\AMCnoCompleteMulti
\begin{questionmult}{3} % тип вопроса (questionmult — множественный выбор) и в фигурных — номер вопроса

Общеизвестно, что потребление мёда Винни-Пухом зависит, при этом положительно, от количества стихов, сочинённых им за день. К сожалению, Винни-Пух забывчив и всегда называет число сочинённых им стихов с ошибкой. Тогда оценка $\beta_1$ в регрессии $Honey_i = \beta_0 + \beta_1 Poems_i + \varepsilon_i $ окажется

%\begin{multicols}{3} % располагаем ответы в 3 колонки
\begin{choices}[o] % опция [o] не рандомизирует порядок ответов
       \correctchoice{Несостоятельной, заниженной}
       \wrongchoice{Несостоятельной, завышенной}
       \wrongchoice{Несостоятельной}
       \wrongchoice{Смещённой, но состоятельной}
       \wrongchoice{Несмещенной, но не состоятельной}
    \end{choices}
   %\end{multicols}
\end{questionmult}
}


% bad style
\element{2017_fall_retake_1}{ % в фигурных скобках название группы вопросов
%  %\AMCnoCompleteMulti
\begin{questionmult}{4} % тип вопроса (questionmult — множественный выбор) и в фигурных — номер вопроса

Из откровений внеземного разума известно, что эндогенности в модели $Y_i = \beta_0 + \beta_1 X_i + \varepsilon_i$ нет. Однако Вовочка нашёл хороший инструмент $z_i$, отвечающий всем требованиям, предъявляемым к инструментам, и оценил $\beta_1$ методом инструментальных переменных. Его оценка $\beta_1$ окажется

\begin{multicols}{3} % располагаем ответы в 3 колонки
\begin{choices} % опция [o] не рандомизирует порядок ответов
       \correctchoice{состоятельной, но не эффективной}
       \wrongchoice{состоятельной и эффективной}
       \wrongchoice{несостоятельной}
       \wrongchoice{состоятельной}
       \wrongchoice{невозможно сказать по имеющимся данным}
    \end{choices}
   \end{multicols}
\end{questionmult}
}





\element{2017_fall_retake_1}{ % в фигурных скобках название группы вопросов
%  %\AMCnoCompleteMulti
\begin{questionmult}{5} % тип вопроса (questionmult — множественный выбор) и в фигурных — номер вопроса

Рассмотрим процесс $Y_t = -0.2 Y_{t-1} + \varepsilon_t$. 5-ое значение автокорреляционной функции равно

\begin{multicols}{1} % располагаем ответы в 3 колонки
\begin{choices} % опция [o] не рандомизирует порядок ответов
       \correctchoice{ -0.00032 }
       \wrongchoice{ 0.00032 }
       \wrongchoice{ 0.2 }
       \wrongchoice{ -0.2 }
       \wrongchoice{ 0 }
    \end{choices}
   \end{multicols}
\end{questionmult}
}



\element{2017_fall_retake_1}{ % в фигурных скобках название группы вопросов
%  %\AMCnoCompleteMulti
\begin{questionmult}{6} % тип вопроса (questionmult — множественный выбор) и в фигурных — номер вопроса

Модель коррекции ошибками имеет следующий вид


\begin{multicols}{2} % располагаем ответы в 3 колонки
\begin{choices} % опция [o] не рандомизирует порядок ответов
       \correctchoice{$\Delta Y_t = \delta + \phi \Delta X_{t-1} - \gamma (Y_{t-1} - \alpha - \beta X_{t-1}) + \varepsilon_t$}
       \wrongchoice{$Y_t = \delta + \phi \Delta X_{t-1} - \gamma (Y_{t-1} - \alpha - \beta X_{t-1}) + \varepsilon_t$}
       \wrongchoice{$\Delta Y_t = \delta - \gamma (Y_{t-1} - \alpha - \beta X_{t-1}) + \varepsilon_t$}
       \wrongchoice{$Y_t = \delta - \gamma (Y_{t-1} - \alpha - \beta X_{t-1}) + \varepsilon_t$}
       \wrongchoice{$\Delta Y_t = \delta + \phi \Delta X_{t-1} - \gamma (Y_{t-1} ) + \varepsilon_t$}
    \end{choices}
   \end{multicols}
\end{questionmult}
}







\element{2017_fall_retake_1}{ % в фигурных скобках название группы вопросов
% \AMCnoCompleteMulti
\begin{questionmult}{7} % тип вопроса (questionmult — множественный выбор) и в фигурных — номер вопроса

Пусть $\varepsilon_t$ - белый шум. Тогда стационарным будет следующий процесс

\begin{multicols}{3} % располагаем ответы в 3 колонки
\begin{choices} % опция [o] не рандомизирует порядок ответов
       \correctchoice{$Y_t = \sum_{i = 0}^{10} \varepsilon_{t-i}$}
       \wrongchoice{$Y_t = 2018t + \varepsilon_t$}
       \wrongchoice{$Y_t = t \varepsilon_t$}
       \wrongchoice{$Y_t = Y_{t-1} - \varepsilon_t$}
       \wrongchoice{$Y_t = 2Y_{t-1} - \varepsilon_t$}
    \end{choices}
   \end{multicols}
\end{questionmult}
}





\element{2017_fall_retake_1}{ % в фигурных скобках название группы вопросов
%  %\AMCnoCompleteMulti
\begin{questionmult}{8} % тип вопроса (questionmult — множественный выбор) и в фигурных — номер вопроса

Процесс случайного блуждания с дрейфом описывается уравнением

\begin{multicols}{2} % располагаем ответы в 3 колонки
\begin{choices} % опция [o] не рандомизирует порядок ответов
       \wrongchoice{$X_t = \mu + X_{t-1} + \varepsilon_t$}
       \wrongchoice{$X_t = \mu + 0.7 X_{t-1} + \varepsilon_t$}
       \wrongchoice{$X_t = X_{t-1} + \varepsilon_t$}
       \wrongchoice{$X_t = 0.7 X_{t-1} + \varepsilon_t$}
       \wrongchoice{$X_t = \mu + \varepsilon_t$}
    \end{choices}
   \end{multicols}
\end{questionmult}
}







\element{2017_fall_retake_1}{ % в фигурных скобках название группы вопросов
 %\AMCnoCompleteMulti
\begin{questionmult}{9} % тип вопроса (questionmult — множественный выбор) и в фигурных — номер вопроса

Если процесс является стационарным в широком смысле, то

\begin{multicols}{3} % располагаем ответы в 3 колонки
\begin{choices} % опция [o] не рандомизирует порядок ответов

       \wrongchoice{ Он является стационарным в узком смысле }
       \wrongchoice{ Для него выполняется основная гипотеза в тесте Дикки-Фуллера }
       \wrongchoice{ Его автоковариационная функция является постоянной }
       \wrongchoice{ Это белый шум }
       \wrongchoice{ Это AR процесс с корнями характеристического уравнения, меньшими 1 }
    \end{choices}
   \end{multicols}
\end{questionmult}
}




\element{2017_fall_retake_1}{ % в фигурных скобках название группы вопросов
 %\AMCnoCompleteMulti
\begin{questionmult}{10} % тип вопроса (questionmult — множественный выбор) и в фигурных — номер вопроса

При оценивании регрессионной модели $Y_t = a_0 + \sum_{j=1}^3 a_j X_{jt} + \varepsilon_t$ по 20 наблюдениям получено значение статистики Дарбина-Уотсона $d = 0.8$. При уровне значимости 1\% это свидетельствует о


\begin{multicols}{3} % располагаем ответы в 3 колонки
\begin{choices} % опция [o] не рандомизирует порядок ответов
       \correctchoice{Попадании в зону неопределенности}
       \wrongchoice{Положительной автокорреляции}
       \wrongchoice{Отрицательной автокорреляции}
       \wrongchoice{Отсутствии автокорреляции}
       \wrongchoice{Тест Дарбина-Уотсона вообще не проверяет наличие автокорреляции}
%       \wrongchoice{3 и 4}
    \end{choices}
   \end{multicols}
\end{questionmult}
}


\AMCnumero{1} % начинаем нумерацию вопросов с 1го

\cleargroup{all}
\copygroup[10]{2017_fall_retake_1}{all}
\insertgroup{all}

\subsubsection{Задачи}


\begin{enumerate}
\item Винни-Пух и Пятачок хотят оценить неизвестный параметр $a$ обобщённым методом моментов.
Винни-Пух наблюдает независимые и одинаково распределённые величины $X_i$
с математическим ожиданием $\E(X_i)=a+2$. А Пяточку известны независимые и
одинаково распределённые величины $Y_i$ с ожиданием $\E(Y_i)=a-1$.

По выборке из 100 величин $X_i$ и из 100 величин $Y_i$ оказалось,
что $\sum X_i = 500$ и $\sum Y_i = -50$.

\begin{enumerate}
\item Найдите оценку обобщённого метода моментов для единичной взвешивающей матрицы.
\item Оцените оптимальную взвешивающую матрицу, если дополнительно известно,
что $\Var(X_i)=a^2 + 25$, $\Var(Y_i)=9$, $\Cov(X_i, Y_i)=-4$.
\end{enumerate}

\item Кролик считает, что процесс $Y_i$ подчиняется уравнению:
\[
Y_t = 7 - 0.9 Y_{t-1} + u_t + u_{t-1},
\]
где процесс $u_t$ — белый шум c дисперсией $\Var(u_t)=\sigma^2_u$.

\begin{enumerate}
\item Есть ли у этого уравнения стационарное решение (является ли данный процесс стационарным)?
Если да, то найдите для него $\E(Y_t)$ и $\Var(Y_t)$.
%\item Приведите пример нестационарного решения данного уравнения.
\item Постройте 95\%-ый предиктивный интервал для $Y_{102}$, если дополнительно известно,
что $Y_{100}=3$, $u_{100}=-1$, а величины $u_t$ имеют нормальное распределение $\cN(0;16)$.
\end{enumerate}

\item Сова пытается \textit{исследовать} зависимость времени,
проведённого Винни-Пухом \textit{на выходе} из дома Кролика,
от количества съеденных блинов с мёдом и блинов со сгущёнкой.
Она оценила модель $vremya_i = \beta_0 + \beta_1 myod_i + \beta_2 sguschenka_i + u_i$
по всем дням, а также отдельно по выходным и будням.

    \begin{center}
 			\begin{tabular}{c|ccc}
				\toprule
				& Будни & Выходные & Вся выборка \\
				\midrule
                RSS & 500 & 400 & 1200 \\
                Наблюдений & 100 & 60 & 160 \\
				\bottomrule
			\end{tabular}
		\end{center}

\begin{enumerate}
\item На уровне значимости 5\% проверьте гипотезу о том, что зависимость не зависит от дня недели.
\item Является ли использованный способ устойчивым к гетероскедастичности?
Если нет, то пошагово опишите способ, устойчивый к гетероскедастичности.
\end{enumerate}


\item
Наблюдения представляют собой случайную выборку. Зависимые переменные $y_{t1}$ и $y_{t2}$ находятся из системы:

\[
\begin{cases}
y_{t1} = \beta_{11} + \beta_{12} x_t + \e_{t1} \\
y_{t2} = \beta_{21} + \beta_{22} z_t + \beta_{23} y_{t1} + \e_{t2}
\end{cases},
\]
где вектор ошибок $\e_t$ имеет совместное нормальное распределение
\[
\e_t \sim \cN\left(
\begin{pmatrix}
  0 \\
  0
\end{pmatrix};
\begin{pmatrix}
  1 & \rho \\
  \rho & 1
\end{pmatrix}
\right)
\]

Эконометресса Анжела оценивает с помощью МНК первое уравнение, а эконометресса Эвридика — второе.
\begin{enumerate}
\item Найдите пределы по вероятности получаемых ими оценок.
\item Будут ли оценки состоятельными?
\end{enumerate}

\item
Исследовательница Мишель строит оценку $\hb_{IV}$ в регрессии $y$ на $x$ с инструментом $z$.
Исследовательница Аграфена строит обычную МНК оценку в регрессии $\hat y = \hb_{x} x + \hb_{w} w$.

Выразите $w$ через $x$, $z$ и $y$ так, чтобы оценка $\hb_{IV}$ Мишель и оценка $\hb_{x}$ Аграфены совпали.



% Каково истинное значение склонности к потреблению, котороы Вы бы ожидали увидеть в этой модели?

\end{enumerate}


\end{document}
